Our first main finding is that LLMs can detect collusive communication. As explained in the methodology section, we made \data{summary_stats/dataset_overview/total_transcripts_int} queries in our large-scale analysis. These queries flagged \data{summary_stats/tagging_performance/llm_tagged_count_int} transcripts. Subsequent validation, to reduce false positives, left us with \data{summary_stats/tagging_performance/llm_validation_tagged_int} transcripts with validated evidence of collusion.

The audit in the top collusive transcripts revealed that \data{summary_stats/tagging_performance/human_audit_false_positive_rate_pct_percentage} are false positives, and approximately \data{summary_stats/tagging_performance/human_audit_rate_pct_percentage} are true positives, according to our screening criteria and human raters. The examples from the introduction show that some of these examples are the subject of ongoing antitrust investigations, lawsuits, or had been considered collusive in previous research \parencite{aryal2022coordinated, harrington2022collusion}. Moreover, \textcite{aryal2022coordinated} give evidence that such communications led to lower capacity and higher prices. We turn to another case study below that illustrates vividly how the LLMs can successfully identify collusive communication.

One way to measure the effectiveness of the LLMs is to count how many of the collusive communications in our benchmark human-rated dataset are rediscovered by the LLMs. In our benchmark dataset, we have \data{summary_stats/tagging_performance/human_benchmark_tagged_count_int} collusive communications that human experts rated as collusive. Out of these, \data{summary_stats/tagging_performance/benchmark_collusive_flagged_count_int} were flagged in the large-scale LLM queries, and \data{summary_stats/tagging_performance/benchmark_collusive_validated_count_int} were validated as collusive.

If we count only the validated transcripts, this means that the LLMs rediscovered \data{summary_stats/tagging_performance/benchmark_collusive_validated_pct_percentage} of the collusive communications in our benchmark dataset. This is, of course, a far shot from 100\%, and illustrates the limitations of the LLMs. However, it is also far above zero, and shows that LLMs have considerable utility in this task.

It is interesting to compare what happened with the transcripts from our benchmark dataset that were flagged as not collusive by humans. Out of these transcripts, only \data{summary_stats/tagging_performance/benchmark_not_collusive_validated_count_int} was validated as collusive, which is \data{summary_stats/tagging_performance/benchmark_not_collusive_validated_pct_percentage} of the total. This is in stark contrast with the benchmark transcripts that humans rated as collusive, which were \data{summary_stats/tagging_performance/benchmark_validation_odds_ratio_int} more likely to be validated as collusive by the LLM.